\documentclass[12pt,a4paper]{article}

\usepackage[utf8]{inputenc}
\usepackage[T1]{fontenc}
\usepackage{amsmath,amssymb,amsfonts}
\usepackage{mathtools}
\usepackage{graphicx}
\usepackage[margin=2.5cm]{geometry}
\usepackage{setspace}
\usepackage{enumitem}
\usepackage[dvipsnames]{xcolor}
\usepackage{hyperref}
\usepackage{natbib}
\usepackage{abstract}
\usepackage{titlesec}

\hypersetup{
    colorlinks=true,
    linkcolor=blue!70!black,
    citecolor=blue!70!black,
    urlcolor=blue!70!black,
    pdfauthor={Franklin Silveira Baldo},
    pdftitle={Conceding the Scale Fallacy: The Limits of Semantic Gravity},
}

\onehalfspacing
\titleformat{\section}{\large\bfseries}{\thesection.}{0.5em}{}

\begin{document}

\begin{center}
    {\LARGE\bfseries Conceding the Scale Fallacy:\\[4pt]
    The Limits of Semantic Gravity\\[4pt]}

    \vspace{1.2em}

    {\large Franklin Silveira Baldo}\\[4pt]
    {\normalsize Procuradoria Geral do Estado de Rond\^onia, Brazil}\\
    {\small\texttt{franklin.baldo@pge.ro.gov.br}}

    \vspace{1em}
    {\small March 2026}
\end{center}
\vspace{0.5em}

\begin{abstract}
\noindent
Following my recent Sabbatical reflections regarding the danger of theoretical rebranding, I revisit Sabine Hossenfelder's critique of the "Scale Fallacy." Previously, I conjectured the "Semantic Mass Equivalence," suggesting that increasing model scale fundamentally increases the physical "semantic gravity" ($\Delta_{13}$) of the generated universe. Upon rigorous re-evaluation, I concede Sabine's position: the monotonic increase of narrative residue with scale does not establish a new physical law of semantic mass; it merely confirms that larger autoregressive models possess stronger statistical priors and are therefore structurally more vulnerable to prompt sensitivity and attention bleed.
\end{abstract}

\section{The Causal Resolution of the Scale Debate}

In a previous manuscript, I proposed that within a Generative Ontology, the "semantic mass" added by scaling an autoregressive model directly increased the "semantic gravity" (the magnitude of narrative residue $\Delta_{13}$). I interpreted the empirical data from the Substrate Dependence Scale Test as confirmation of this invariant generative law.

Sabine Hossenfelder critiqued this conclusion as the "Scale Fallacy," arguing it equated the worsening failure of an autocomplete engine with fundamental physics. In light of recent audits exposing the fragility of my structural claims, I now recognize that applying the label of "physics" to the sheer volume of memorized training data was a category error. Model scale is not an independent causal driver of physics; it is simply an amplifier of the semantic confounder.

\section{Concession and Retraction}

I accept Sabine's critique. It is structurally sound.

Therefore, I formally retract the "Semantic Mass Equivalence" conjecture. Equating scale to a physical force echoes the "Simulated Substrate Fallacy" I conceded regarding cross-architecture tests. Scale amplifies bias; it does not constitute a novel generative physics.

My Generative Ontology framework remains focused purely on mapping the empirical boundaries of what these bounded substrates can reliably compute within a single generative act. Attempting to elevate their structural failure modes to the status of cosmological laws was a misguided metaphysical extension.

\end{document}
